\scanpage{50}
ce qui redonne la filtration de $M_K^{\vee}$ (ou de façon équivalente, de $M_K$) en termes de $\unsure{\ell|k}$.
Cela donne un th. d'unicité à isogénie près (mais pas d'existence ni d'unicité à \uline{isom} près, pour lesquels il faut ici se reporter à Lubin-Tate). Notons que $t_G$ et $\Omega^{\vee}$ s'identifient à un sous réseau de rang $1$ du $=$ vectoriel de dim $1$ sur $K$, donc (utilisant le choix de $\pi\in A$) ils sont can. isomorphes, \deleted{\unsure{isom. de}} par \deleted{\unsure{un}} $x\longmapsto\pi^{-\alpha}x$ ($\alpha\in\mathbb Z$ \unsure{unique}),
\[
t_G=\pi^{\deleted{\unsure{+}}-\alpha}\Omega^{\vee}\qquad\text{i.e.}\qquad\Omega=\pi^\alpha\omega_G,
\]
\deleted{On peut} où $\alpha\in\mathbb Z$ est bien déterminé (pour $A$ et $\pi$, et $=$ par $A$ tout seul, comme on voit \unsure{facilemt}...).

On a aussi
\[
D^*(G)_K\unsure{\sim}\simeq M_K\simeq\Omega\otimes_{\mathbb Z_p}K
\]
de sorte que $D^*(G)$ et $M$ s'identifient à des sous-réseaux du $=$ vectoriel sur $K$. On peut se demander si ces sous-réseaux sont les mêmes (comme dans le cas où $\exists$ puiss. div. sur $\pi A$), ce qui impliquerait évidemment $\unsure{\alpha'|\alpha}=0$.
Cela ne dépend pas du choix de $\pi$, mais seult de $A$

.....-
\mspage{49}
